% !TeX encoding = UTF-8
% Français de France. Charger un seul archétype depuis le document principal.
\newcommand{\cvsummary}{Ingénieur en mécanique et mécatronique associant conception CAO avancée, analyse aérodynamique et réalisation de prototypes. Expérience en conception mécanique, fabrication, intelligence artificielle et ingénierie de projets industriels, avec une pratique de la commande embarquée et des systèmes robotiques pour les intégrer à la mécanique.}

% --- COMPÉTENCES ---
\newcommand{\skills}{
    \section*{COMPÉTENCES TECHNIQUES}
    {\small
    \textbf{Conception mécanique et CAO :} SolidWorks (CSWA), Solid Edge, CATIA V5, AutoCAD, Siemens NX, Creo, Fusion 360, modélisation paramétrique 3D, plans de fabrication 2D, cotation dimensionnelle et géométrique (GD\&T, ISO), interfaces mécaniques, conception pour injection plastique.

    \vspace{1pt}
    \textbf{Simulation structurelle et dynamique :} ANSYS Mechanical/Fluent, SolidWorks Simulation, COMSOL, analyse par éléments finis (FEA), mécanique des fluides numérique (CFD), analyses structurelles et vibratoires, analyse de maillage, cinématique, MATLAB/Simulink, analyse aérodynamique.

    \vspace{1pt}
    %\textbf{Intégration de systèmes et prototypage :} Assemblages électromécaniques, systèmes aérospatiaux, structures d'aéronefs, capteurs et actionneurs embarqués, intégration mécanique-électrique, assemblage de prototypes, microcontrôleurs, MEMS.

    \vspace{1pt}
    \textbf{Fabrication et matériaux :} Usinage CNC, fabrication additive (PLA/TPU), moulage par injection, composites en fibre de carbone, tôlerie, sélection de matériaux, conception pour la fabrication et l'assemblage (DFM/DFA), métrologie, fabrication et assemblage de prototypes.

    \vspace{1pt}
    \textbf{Programmation et IA :} Python, MATLAB, C++, SQL, JavaScript, Java, HTML ; analyse de données, automatisation de processus, apprentissage automatique, classification, réseaux de neurones, architectures Transformer, processus gaussiens.

    \vspace{1pt}
    \textbf{Essais et documentation technique :} Essais de chute et de vibrations, validation expérimentale, étalonnage de capteurs, plans techniques, rapports d'ingénierie, LaTeX, MS Project, Jira, Notion, Six Sigma White Belt.

    \vspace{1pt}
    \textbf{Langues :} Espagnol langue maternelle, anglais C2, français C2, italien A1, allemand A1.
    }
}


% --- PROJETS ---
% --- PROJETS ---
\newcommand{\vawttitle}{VAWT -- Récupération d'énergie éolienne urbaine}
\newcommand{\vawtdates}{09/2024 - 05/2025}
\newcommand{\vawtlocation}{Abou Dabi, ÉAU}
\newcommand{\vawtdescOne} {Conception et optimisation CFD d'une éolienne à axe vertical hybride à portance et traînée, avec pales hélicoïdales de 60° pour la récupération du vent urbain (2,25--24 m/s) ; fabrication de pales en fibre de carbone multipliant la tension moyenne par 3,2 par rapport au PLA et atteignant 13,02 V en crête. Projet de fin d'études avec brevet en instance.}
\newcommand{\vawtdescTwo}{Développement d'une éolienne à axe vertical hybride à portance et traînée pour exploiter les vents urbains, avec optimisation de géométries hélicoïdales spécifiques de 60° par CAO et analyses CFD de portance, traînée et couple sur des vitesses de vent de 2,25--24 m/s. Développement de pales composites en fibre de carbone multipliant la tension moyenne par 3,2 et atteignant 13,02 V en crête par rapport au prototype initial en PLA.}

\newcommand{\fsaetitle}{Formula SAE -- Aérodynamique active et conception d'un châssis en fibre de carbone}
\newcommand{\fsaedates}{Automne 2024 -- Aujourd'hui}
\newcommand{\fsaelocation}{Abou Dabi, ÉAU}
\newcommand{\fsaedescOne}{Conception des géométries aérodynamiques et des composants de châssis en fibre de carbone sous SolidWorks, avec validation structurelle par éléments finis sous ANSYS. Optimisation de l'aérodynamique des ailerons sous ANSYS Fluent, atteignant 230 N d'appui ; développement actuel de l'aérodynamique active pour UniOvi Racing.}
\newcommand{\fsaedescTwo}{Évaluation de la cinématique de suspension et coordination des interfaces structurelles avec les sous-équipes aérodynamique et fabrication, en reliant exigences de mouvement du véhicule et intégration des composants.}

\newcommand{\tactitle}{Mini-réfrigérateur -- Système thermoélectrique TEC à deux étages}
\newcommand{\tacdates}{Printemps 2025}
\newcommand{\taclocation}{Abou Dabi, ÉAU}
\newcommand{\tacdescOne}{Conception d'un mini-réfrigérateur sans fluide frigorigène autour d'un module TEC à deux étages, avec dissipateurs imprimés en 3D et structure de support pour intégrer les composants thermiques.}
\newcommand{\tacdescTwo}{Évaluation en temps réel des performances de refroidissement grâce à l'instrumentation thermique intégrée au prototype, en reliant les choix de conception mécanique au comportement observé.}

\newcommand{\orbitaltitle}{OrbitalGrid -- Système SIG pour applications spatiales}
\newcommand{\orbitaldates}{06/2024 - Aujourd'hui}
\newcommand{\orbitallocation}{Abou Dabi, ÉAU}
\newcommand{\orbitaldescOne}{Développement d'un système cartographique fondé sur SQL combinant numérisation 3D et données satellitaires, en réunissant géométrie physique et informations géospatiales dans un processus commun de données d'ingénierie.}
\newcommand{\orbitaldescTwo}{Organisation d'informations cartographiques sous PostgreSQL par ingénierie des données spatiales, pour intégrer géométries numérisées et observations satellitaires.}

\newcommand{\softrobtitle}{Exosquelette robotique pour la rééducation après un AVC}
\newcommand{\softrobdates}{04/2024 - 05/2024}
\newcommand{\softroblocation}{Abou Dabi, ÉAU}
\newcommand{\softrobdescOne}{Conception d'un exosquelette de doigts inspiré de l'origami, avec actionnement par fils SMA et polymère Dragon Skin, combinant géométrie souple et éléments actifs pour la rééducation.}
\newcommand{\softrobdescTwo}{Évaluation de la stabilité et de la durabilité du mécanisme par simulations sous SolidWorks, avant de poursuivre la conception mécanique.}

\newcommand{\exoplanettitle}{Simulateur de transits d'exoplanètes}
\newcommand{\exoplanetdates}{Printemps 2025}
\newcommand{\exoplanetlocation}{Abou Dabi, ÉAU}
\newcommand{\exoplanetdescOne}{Développement d'un simulateur Python paramétrable générant des courbes de lumière de transits à partir de données de mécanique orbitale et de physique stellaire, en transformant des relations physiques en résultats numériques.}
\newcommand{\exoplanetdescTwo}{Comparaison de différentes configurations étoile--planète grâce à des paramètres de modèle accessibles et à la représentation des réponses temporelles sous Matplotlib.}

\newcommand{\windTurbinetitle}{Éolienne flottante en mer -- Stabilisation active}
\newcommand{\windTurbinedates}{05/2026 - 06/2026}
\newcommand{\windTurbinelocation}{Besançon, France}
\newcommand{\windTurbinedescOne}{Conception et réalisation sous SolidWorks d'un prototype d'éolienne semi-submersible à trois colonnes à l'échelle 1:1000, avec stabilité passive en tangage et roulis, GM = 5,38 cm et raideur de rappel K = 0,0607 N$\cdot$m/rad ; intégration d'un MPU6050 et d'un cardan actionné par servomoteurs pour stabiliser activement la nacelle lors d'essais en bassin à houle.}

\newcommand{\carbonmetertitle}{CARBONMETRE -- Écoconception et mesure carbone}
\newcommand{\carbonmeterdates}{09/2025 - 12/2025}
\newcommand{\carbonmeterlocation}{Besançon, France}
\newcommand{\carbonmeterdescOne}{Association de la fabrication, de l'écoconception et de la métrologie dans un projet de mesure carbone, en reliant méthodes de fabrication, mesure mécanique et évaluation environnementale.}
\newcommand{\carbonmeterdescTwo}{Application des principes de fabrication industrielle au projet CARBONMETRE, en intégrant les considérations de procédé aux travaux de mesure et de conception durable.}

\newcommand{\hvactitle}{Amélioration thermique d'une chambre en résidence -- Refroidissement passif et évaporatif}
\newcommand{\hvacdates}{06/2026}
\newcommand{\hvaclocation}{Besançon, France}
\newcommand{\hvacdescOne}{Réduction de la température d'une chambre de 9 m$^2$ jusqu'à 3 \textdegree{}C par rapport à la situation avant intervention, par conception et fabrication d'un écran solaire extérieur et d'un refroidisseur à eau gelée entraîné par ventilateur.}
\newcommand{\hvacdescTwo}{Blocage du rayonnement solaire direct entre 14 h et 21 h par papier aluminium réfléchissant, PSE de 1 cm et carton ondulé ; fabrication d'un caisson de refroidissement à sept sorties pour répartir le flux d'air sur une source froide d'eau gelée.}

\newcommand{\controlBoxtitle}{Installation électrique complète -- Équipements pour événements de taille moyenne}
\newcommand{\controlBoxdates}{06/2023 - 08/2023}
\newcommand{\controlBoxlocation}{Mexico, Mexique}
\newcommand{\controlBoxdescOne}{Réduction de 30\% du temps de montage du système électrique par rapport à l'installation précédente, grâce à la conception d'un coffret de commande de 12 entrées pour des installations événementielles de taille moyenne.}
\newcommand{\controlBoxdescTwo}{Intégration des commandes électriques à automate programmable à l'ensemble audiovisuel, en traduisant les exigences de fonctionnement d'événements en direct en matériels et en exploitation coordonnée du système.}

\newcommand{\hyperlooptitle}{NYC Hyperloop -- Prototype à échelle réduite}
\newcommand{\hyperloopdates}{01/2024 - 05/2024}
\newcommand{\hyperlooplocation}{New York, États-Unis}
\newcommand{\hyperloopdescOne}{Pilotage de l'équipe mécanique d'un prototype Hyperloop à échelle réduite pour un concept de liaison New York--Long Island, en coordonnant la conception du véhicule autour d'une suspension électrodynamique Inductrack.}
\newcommand{\hyperloopdescTwo}{Pilotage du développement mécanique d'un système Hyperloop à échelle réduite pour Manhattan--Long Island, avec conception d'une capsule à lévitation Inductrack par réseaux de Halbach, roues auxiliaires et propulsion/freinage électromagnétiques. Évaluation de trois architectures par CAO et simulation, avec contribution au prototypage CNC et impression 3D et à l'intégration du système.}

\newcommand{\dronetitle}{Quadricoptère fonctionnel -- Structure allégée}
\newcommand{\dronedates}{03/2024 - 05/2024}
\newcommand{\dronelocation}{New York, États-Unis}
\newcommand{\dronedescOne}{Conception et impression 3D d'un quadricoptère à quatre bras en PLA et TPU, avec analyse modale par éléments finis et essais de balayage sinusoïdal, de vibrations aléatoires et de chute. Validation de l'intégrité structurelle avec une fréquence propre de 23,27 Hz et tenue à un essai de vibrations aléatoires de 10 min à 0,04 g$^2$/Hz entre 5 et 200 Hz.}
\newcommand{\dronedescTwo}{Évaluation de la robustesse de la structure par essais de chute et de vibrations selon les critères du projet, complétée par des analyses structurelles et aérodynamiques sous SolidWorks et ANSYS.}

\newcommand{\trentxwbtitle}{Rolls-Royce Trent XWB -- Modèle 3D fonctionnel}
\newcommand{\trentxwbdates}{01/2024 - 05/2024}
\newcommand{\trentxwblocation}{New York, États-Unis}
\newcommand{\trentxwbdescOne}{Codéveloppement d'un assemblage détaillé du turboréacteur à double flux Trent XWB sous SolidWorks, avec modélisation des composants de la soufflante d'entrée au cœur et à l'échappement. Simulations structurelles et aérodynamiques sous ANSYS pour évaluer les contraintes de l'arbre et les performances de la soufflante, avec production de plans techniques et de rendus.}
\newcommand{\trentxwbdescTwo}{Extension de l'évaluation du moteur au-delà de la géométrie grâce à des simulations aérodynamiques sous ANSYS, en combinant modélisation d'assemblage et analyse des écoulements.}

\newcommand{\liostitle}{LIOS -- Assistant vocal avec IA}
\newcommand{\liosdates}{12/2024 - Aujourd'hui}
\newcommand{\lioslocation}{Abou Dabi, ÉAU}
\newcommand{\liosdescOne}{Développement d'une application d'assistant vocal en Python intégrant traitement automatique du langage, données en temps réel et API externes, en passant de fonctions isolées à un système interactif.}
\newcommand{\liosdescTwo}{Extension de l'assistant au-delà des requêtes isolées par mémoire contextuelle, services météorologiques et tâches automatisées dans une architecture Python modulaire.}

\newcommand{\fitpaltitle}{FitPal -- Dispositif intelligent d'entraînement physique}
\newcommand{\fitpaldates}{09/2021 - 12/2021}
\newcommand{\fitpallocation}{Abou Dabi, ÉAU}
\newcommand{\fitpaldescOne}{Codéveloppement d'un dispositif d'entraînement physique sur M5Stack pour analyser les exercices en temps réel, en pilotant la conception du produit durant son développement et en intégrant matériel embarqué et mesures destinées à l'utilisateur.}
\newcommand{\fitpaldescTwo}{Connexion du dispositif physique à une interface mobile via Arduino IoT App, avec intégration du transfert de données et de l'accès à distance au produit.}

\newcommand{\covidtitle}{Cabine de désinfection automatisée}
\newcommand{\coviddates}{09/2020 - 02/2021}
\newcommand{\covidlocation}{Mexico, Mexique}
\newcommand{\coviddescOne}{Conception d'un prototype de cabine de désinfection automatisée associant commande Arduino et matériel de brumisation, en transformant un besoin fonctionnel en système physique.}
\newcommand{\coviddescTwo}{Intégration de la communication sans fil aux commandes de brumisation de la cabine, en reliant les équipements d'acheminement du fluide et la commande électronique dans le prototype.}

\newcommand{\rovertitle}{Rover téléopéré -- Numérisation spatiale 3D}
\newcommand{\roverdates}{09/2026 - Aujourd'hui}
\newcommand{\roverlocation}{Oviedo, Espagne}
\newcommand{\roverdescOne}{Conception d'un châssis de rover en ABS pour moulage par injection, intégrant un scanner Flexx2, M5Stack, Raspberry Pi et une chaîne de traction commandée par FPGA, avec optimisation des interfaces mécaniques et de l'assemblage. Intégration de la numérisation 3D, de la commande distante TCP/IP, d'une direction Ackermann et d'un rapport de réduction moteur de 1:20.}
\newcommand{\roverdescTwo}{Intégration en cours de l'acquisition de nuages de points avec visualisation et commande distantes sur TCP/IP, en reliant la plateforme mécanique à la communication M5Stack--Raspberry Pi ; validation de bout en bout encore en cours.}

\newcommand{\microsensortitle}{Microcapteur piézorésistif de pression de détonation}
\newcommand{\microsensordates}{05/2026 - 06/2026}
\newcommand{\microsensorlocation}{Besançon, France}
\newcommand{\microsensordescOne}{Réduction des itérations de géométrie et de masques lithographiques en remplaçant les calculs manuels de dimensionnement du capteur par un outil Python imposant des contraintes de facteur de jauge, $\Delta R/R_0$ et résistivité.}
\newcommand{\microsensordescTwo}{Traduction d'une interface spécifiée à 50 $\Omega$ et d'un objectif de tenue à 10 GPa en choix de matériaux, dessin de masques et gamme de fabrication, en conciliant performances électriques, exigences mécaniques, coût et contraintes de salle blanche.}

\newcommand{\missiletitle}{Simulateur de dynamique d'interception de missiles}
\newcommand{\missilelocation}{Besançon, France}
\newcommand{\missiledescOne}{Développement d'un simulateur 3D d'interception sous MATLAB avec navigation proportionnelle, intégration numérique et détection d'impact, en transformant les équations des systèmes dynamiques en visualisations de trajectoires pour évaluer des scénarios d'interception.}

% --- EXPÉRIENCE ---
\newcommand{\carbonorg}{Energy Innovation Capital}
\newcommand{\carbontitle}{Analyste de marché -- Captage du carbone}
\newcommand{\carbondates}{10/2023 - 01/2024}
\newcommand{\carbonlocation}{Abou Dabi, ÉAU}
\newcommand{\carbonbullets}{%
  \item Élaboration de recommandations d'investissement sur des jeunes entreprises de captage du carbone par recherche technique et analyse de marché pour un fonds de capital-risque d'Abou Dabi.
  \item Transformation des avancées scientifiques et technologiques en présentations pour les analystes seniors, afin de rendre les résultats d'ingénierie utiles aux discussions d'investissement.
  \item Évaluation des tendances de marché, des moteurs technologiques et des risques d'investissement pour replacer les solutions CCUS émergentes dans leur contexte commercial.
  \item Mise en relation des dimensions techniques et commerciales des jeunes entreprises, par synthèse des résultats de recherche en recommandations stratégiques.
}

\newcommand{\spacelaborg}{NYU Abu Dhabi, Centre des sciences spatiales}
\newcommand{\spacelabtitle}{Ingénieur et chercheur R\&D -- Plateforme de sciences planétaires}
\newcommand{\spacelabdates}{05/2025 - 09/2025}
\newcommand{\spacelablocation}{Abou Dabi, ÉAU}
\newcommand{\spacelabbullets}{%
  \item Déploiement d'un atlas interactif de Mars fondé sur des images multispectrales réelles d'Emirates Mars Mission HOPE, avec développement de l'\href{https://romvio.github.io/mars-viewer/}{application web et mobile} en React et JavaScript et soumission à l'App Store d'Apple.
  \item Amélioration de 40\% du traitement des images planétaires par compression optimisée et correction des pixels assistée par apprentissage automatique, en assemblant les images HOPE par calibration des recouvrements, mosaïquage et normalisation de la luminosité.
  %\item Collecte d'échantillons de sols et de sédiments sur des sites désertiques des ÉAU par carottage au tube et excavation, avec documentation du terrain et des sédiments pour rechercher de possibles analogues du régolithe martien.
  %\item Déploiement d'un atlas interactif de Mars fondé sur des images multispectrales réelles d'Emirates Mars Mission HOPE, avec développement de l'\href{https://romvio.github.io/mars-viewer/}{application web et mobile} en React et JavaScript et soumission à l'App Store d'Apple.
  %\item Amélioration de 40\% de l'efficacité du traitement d'images par compression optimisée et correction assistée par apprentissage automatique, avec calibration des recouvrements, mosaïquage et normalisation de la luminosité.
  %\item Collecte d'échantillons de sols et de sédiments sur des sites désertiques des ÉAU par carottage au tube et excavation pour rechercher de possibles analogues du régolithe martien.
}

\newcommand{\geoorg}{Global Education NYUAD}
\newcommand{\geotitle}{Chargé de coordination Global Education -- Étudiants en mobilité entrante}
\newcommand{\geodates}{12/2025 - 05/2025}
\newcommand{\geolocation}{Abou Dabi, ÉAU}
\newcommand{\geobullets}{%
  \item Organisation de sorties universitaires et culturelles aux ÉAU par coordination logistique pour les étudiants internationaux en mobilité entrante à NYU Abu Dhabi.
  \item Pilotage de la planification et de la répartition des responsabilités de l'équipe de coordinateurs, en reliant les exigences du programme aux opérations pratiques des groupes.
  \item Intégration de protocoles de sécurité à la préparation des sorties, en adaptant les décisions logistiques aux besoins des participants et aux précautions opérationnelles.
  \item Coordination des échanges entre étudiants en mobilité et services universitaires, avec un point de contact direct pour l'accompagnement du programme.
}

\newcommand{\geoOneorg}{Global Education NYUAD}
\newcommand{\geoOnetitle}{Chargé de coordination Global Education -- Étudiants de première année}
\newcommand{\geoOnedates}{09/2026 - 12/2024}
\newcommand{\geoOnelocation}{Abou Dabi, ÉAU}
\newcommand{\geoOnebullets}{%
  \item Animation de séances d'accueil de première année, en organisant les informations universitaires sous forme de présentations structurées de la vie du campus.
  \item Accompagnement de l'intégration étudiante par mentorat et conseil continus entre pairs pour faciliter l'accès aux ressources universitaires.
  \item Conduite de visites du campus présentant les installations et services, avec adaptation des informations pratiques aux besoins des nouveaux étudiants.
  \item Organisation d'événements universitaires et culturels inclusifs, en traduisant les objectifs de cohésion communautaire en activités planifiées.
}

\newcommand{\hcorg}{Centre de santé de NYUAD}
\newcommand{\hctitle}{Emploi étudiant -- Responsable des réseaux sociaux}
\newcommand{\hcdates}{09/2024 - 05/2025}
\newcommand{\hclocation}{Abou Dabi, ÉAU}
\newcommand{\hcbullets}{%
  \item Réalisation de campagnes de santé ciblées par adaptation des informations de bien-être à la communauté étudiante de NYU Abu Dhabi.
  \item Conception et publication de supports numériques de campagne, en transformant les initiatives de santé en contenus accessibles sur les réseaux sociaux.
  \item Gestion des demandes étudiantes comme contact principal du centre de santé, en orientant les questions vers les services compétents.
  \item Contribution à la résolution des demandes par suivi direct, en associant communication et coordination des services.
}

\newcommand{\alpesorg}{Alpes Technologies México}
\newcommand{\alpestitle}{Ingénieur projets junior -- Gestion d'entrepôt et systèmes électriques}
\newcommand{\alpesdates}{06/2024 - 09/2024}
\newcommand{\alpeslocation}{Mexico, Mexique}
\newcommand{\alpesbullets}{%
  \item Développement intégral d'une application de gestion d'entrepôt avec codes QR en HTML/JS et stockage sur les services Google, reliant les flux physiques de stock au suivi numérique en temps réel et améliorant de 60\% l'efficacité opérationnelle de l'entrepôt.
  \item Fabrication et installation de batteries de condensateurs industrielles avec filtres harmoniques après analyse des harmoniques et pics de tension des postes, atteignant un facteur de puissance de 0,99 et réduisant les coûts énergétiques de 15\% et les arrêts électriques de 25\%.
  \item Diagnostic électrique après l'incendie d'un complexe hôtelier par mesures sur site, avec schémas unifilaires et recommandations correctives présentés à la direction générale.
  %\item Production de schémas unifilaires et de recommandations correctives pour un complexe hôtelier sinistré par un incendie, par mesures sur site et présentation des résultats aux directions de l'entreprise et du client.
}

\newcommand{\epigsaorg}{EPIGSA}
\newcommand{\epigsatitle}{Assistant en ingénierie des données -- SIG et hydraulique}
\newcommand{\epigsadates}{05/2024 - 07/2024}
\newcommand{\epigsalocation}{Mexico, Mexique}
\newcommand{\epigsabullets}{%
  \item Réduction de 25\% des erreurs de données d'infrastructures par automatisation des contrôles qualité pour des projets municipaux de traitement des eaux.
  \item Intégration et validation de jeux de données SIG sous PostgreSQL, avec préparation d'informations géospatiales pour l'ingénierie hydraulique et des infrastructures.
  %\item Amélioration de la fiabilité des données d'entrée des projets par vérification des enregistrements d'infrastructures lors de leur intégration en base.
  %\item Mise en relation des informations d'ingénierie et des actifs physiques par intégration de données SIG et validation automatisée.
}

\newcommand{\solarshiporg}{Solar Ship}
\newcommand{\solarshiptitle}{Stagiaire en stratégie et conception de mobilité aérienne}
\newcommand{\solarshipdates}{05/2023 - 08/2023}
\newcommand{\solarshiplocation}{Toronto, Canada}
\newcommand{\solarshipbullets}{%
  \item Développement de concepts d'infrastructures de logistique aérienne utilisant des technologies à base d'hydrogène pour les zones isolées de Zambie et de Tanzanie, en traduisant les besoins d'accès hors réseau en travaux de conception aérospatiale et d'infrastructure.
  \item Contribution à un système aéronautique entièrement électrique pour l'initiative Peace+Freedom, en adaptant le développement du véhicule aux besoins de connectivité des communautés isolées et en traitant la logistique de soutien avec le véhicule aérien.
  %\item Intégration d'objectifs d'ingénierie et de durabilité sociale dans un projet multinational de mobilité pour des communautés hors réseau.
}

\newcommand{\crousorg}{CROUS Besançon Bourgogne Franche-Comté}
\newcommand{\croustitle}{Emploi étudiant -- Restauration universitaire}
\newcommand{\crousdates}{10/2024 - Aujourd'hui}
\newcommand{\crouslocation}{Besançon, France}
\newcommand{\crousbullets}{%
  \item Maintien du rythme en cuisine aux heures de pointe par coordination de l'appui à la préparation des aliments et des tâches internes dans une activité à fort volume.
  \item Respect des pratiques d'hygiène et de sécurité alimentaire par nettoyage, désinfection et plonge en cuisine.
  \item Contribution à un service continu par réapprovisionnement adapté à la demande opérationnelle.
  \item Adaptation des tâches de cuisine et d'appui au service aux priorités changeantes, en contribuant à une exécution coordonnée sous contrainte de temps.
}

\newcommand{\rallyorg}{Championnat du monde des rallyes -- Guanajuato}
\newcommand{\rallytitle}{Commissaire technique officiel FIA}
\newcommand{\rallydates}{11/2019 - 04/2020}
\newcommand{\rallylocation}{Guanajuato, Mexique}
\newcommand{\rallybullets}{%
  \item Marquage officiel des pneumatiques et inspections techniques au Championnat du monde des rallyes FIA à Guanajuato, selon les procédures prescrites de vérification.
  \item Contribution à la conformité des véhicules de compétition par contrôles au regard des exigences de l'épreuve.
  \item Coordination des inspections avec les commissaires responsables et les pilotes sous la pression temporelle d'un week-end de compétition.
  \item Application d'un contrôle qualité procédural par marquage et vérification systématiques des pneumatiques, au service d'opérations techniques réglementées.
}

\newcommand{\papaloteorg}{Papalote Museo del Niño}
\newcommand{\papalotetitle}{Chargé d'accueil et de médiation}
\newcommand{\papalotedates}{05/2019 - 07/2019}
\newcommand{\papalotelocation}{Mexico, Mexique}
\newcommand{\papalotebullets}{%
  \item Présentation d'expositions scientifiques et techniques aux enfants et aux adultes, en transformant les concepts techniques en explications pratiques pour les visiteurs.
  \item Amélioration de l'expérience de visite par propositions et changements mis en œuvre dans les services du musée.
  \item Adaptation des explications aux différents âges et besoins pour faciliter la participation d'un public diversifié aux expositions techniques.
  \item Mise en relation des démonstrations scientifiques et de l'interaction directe avec les visiteurs pour faciliter la compréhension de mécanismes et de concepts inconnus.
}

